% Source: https://artofproblemsolving.com/community/c3h1600643_mathcounts_2014_national_team_4
% Stable identity: aops-topic-1600643-q01
\textbf{Closed form.} The first terms are $0,4,16,48,128,320,768,\dots$. We claim $a_n=(n-1)2^n$. It holds for $n=1,2$ ($0$ and $4$). If it holds for $n-1$ and $n-2$, then
\[
a_n=4\left((n-2)2^{n-1}-(n-3)2^{n-2}\right)=4\cdot2^{n-2}\left(2(n-2)-(n-3)\right)=2^n(n-1),
\]
so it holds for all $n$ by induction.

\textbf{Powers of $2$.} For $n\ge2$, $a_n=(n-1)2^n$ is a power of $2$ exactly when $n-1$ is a power of $2$ (any odd factor of $n-1$ greater than $1$ would divide $a_n$). And $a_1=0$ is not a power of $2$. So $a_n$ is a power of $2$ exactly when $n=2^m+1$ with $m\ge0$.

The largest such $n$ below $100$ is $n=64+1=\boxed{65}$, with $a_{65}=64\cdot2^{65}=2^{71}$.
